\documentclass[a4paper,10pt]{article}
\usepackage[utf8]{inputenc}
\usepackage[T1]{fontenc}
\usepackage[margin=1.8cm,top=2.2cm,bottom=2.2cm]{geometry}
\usepackage{booktabs}
\usepackage{colortbl}
\usepackage{xcolor}
\usepackage{array}
\usepackage{multirow}
\usepackage{microtype}
\usepackage{parskip}
\usepackage{enumitem}
\usepackage{fancyhdr}
\usepackage{titlesec}
\usepackage{tabularx}
\usepackage{amsmath}
\usepackage{lmodern}
\usepackage{hyperref}

% ── Colour palette ──────────────────────────────────────────────────
\definecolor{navy}{RGB}{0,32,96}
\definecolor{teal}{RGB}{0,112,120}
\definecolor{lightgray}{RGB}{245,245,245}
\definecolor{midgray}{RGB}{180,180,180}
\definecolor{alertred}{RGB}{180,0,0}
\definecolor{kpibg}{RGB}{0,32,96}

% ── Section formatting ───────────────────────────────────────────────
\titleformat{\section}{\normalsize\bfseries\color{navy}}{}{0em}{\MakeUppercase}[\vspace{-4pt}\color{midgray}\hrule\vspace{4pt}]
\titleformat{\subsection}{\small\bfseries\color{teal}}{}{0em}{}
\setlist{noitemsep,topsep=2pt,leftmargin=*}
\renewcommand{\arraystretch}{1.3}

% ── Header / footer ─────────────────────────────────────────────────
\pagestyle{fancy}\fancyhf{}
\fancyhead[L]{\small\textcolor{navy}{\textbf{Indonesia Data Centre Sector --- Fiscal Cost of Tax Incentives}}}
\fancyhead[R]{\small\textcolor{midgray}{March 2026 | Confidential}}
\fancyfoot[C]{\small\textcolor{midgray}{\thepage}}
\renewcommand{\headrulewidth}{0.6pt}
\renewcommand{\footrulewidth}{0pt}

\hypersetup{colorlinks=false}

\begin{document}

% ═══════════════════════════════════════════════════════════════════
%  PAGE 1 — FISCAL COST SUMMARY
% ═══════════════════════════════════════════════════════════════════

{\centering
  {\LARGE\bfseries\color{navy} Indonesia Data Centre Sector}\\[3pt]
  {\large\color{teal} Fiscal Cost of Tax Incentives --- Portfolio Summary}\\[2pt]
  {\small\color{midgray} 45 Facilities (23 Operational \textbar{} 12 Under Construction \textbar{} 10 Planned) \textbar{} March 2026}\\[6pt]
}

\noindent\color{navy}\rule{\linewidth}{1.5pt}

\vspace{8pt}

% ── KPI row ─────────────────────────────────────────────────────────
\noindent
\begin{tabularx}{\linewidth}{>{\centering\arraybackslash}X>{\centering\arraybackslash}X>{\centering\arraybackslash}X>{\centering\arraybackslash}X>{\centering\arraybackslash}X}
  \cellcolor{kpibg}\textcolor{white}{\textbf{\large IDR 373.9T}} &
  \cellcolor{kpibg}\textcolor{white}{\textbf{\large IDR 118.2T}} &
  \cellcolor{kpibg}\textcolor{white}{\textbf{\large IDR 19.1T}}  &
  \cellcolor{kpibg}\textcolor{white}{\textbf{\large IDR 0.8T}}   &
  \cellcolor{alertred}\textcolor{white}{\textbf{\large IDR 138.1T}} \\
  \cellcolor{kpibg}\textcolor{white}{\scriptsize Total Investment} &
  \cellcolor{kpibg}\textcolor{white}{\scriptsize Tax Holiday Cost} &
  \cellcolor{kpibg}\textcolor{white}{\scriptsize Super Deduction Cost} &
  \cellcolor{kpibg}\textcolor{white}{\scriptsize SEZ VAT Exemption} &
  \cellcolor{alertred}\textcolor{white}{\scriptsize \textbf{Grand Total Fiscal Cost}} \\
  \cellcolor{kpibg}\textcolor{white}{\scriptsize (USD 24.1B)} &
  \cellcolor{kpibg}\textcolor{white}{\scriptsize (over holiday periods)} &
  \cellcolor{kpibg}\textcolor{white}{\scriptsize (6 yrs post-holiday)} &
  \cellcolor{kpibg}\textcolor{white}{\scriptsize (Nongsa KEK only)} &
  \cellcolor{alertred}\textcolor{white}{\scriptsize (all incentives combined)}\\
\end{tabularx}

\vspace{10pt}

% ── Tier breakdown ───────────────────────────────────────────────────
\section{Fiscal Cost Breakdown by Investment Tier (PMK 69/2024)}

\small
\begin{tabular}{lrrrrrrr}
  \toprule
  \textbf{Investment Tier} & \textbf{N} & \textbf{Investment} & \textbf{Holiday} & \textbf{Holiday} & \textbf{Super Ded.} & \textbf{VAT} & \textbf{Total} \\
  & & \textbf{(IDR T)} & \textbf{Years} & \textbf{Cost (IDR T)} & \textbf{Cost (IDR T)} & \textbf{(IDR T)} & \textbf{(IDR T)} \\
  \midrule
  Below IDR 100B          & 2  &   0.14 &  0  &   0.000 &  0.000 & 0.000 &  0.000 \\
  IDR 100B--500B          & 2  &   0.64 &  5  &   0.043 &  0.042 & 0.018 &  0.103 \\
  IDR 500B--1T            & 3  &   2.34 &  5  &   0.157 &  0.079 & 0.000 &  0.236 \\
  IDR 1T--5T              & 19 &  47.72 &  7  &   6.379 &  2.283 & 0.119 &  8.781 \\
  IDR 5T--15T             & 14 &  98.36 & 10  &  18.460 &  4.616 & 0.659 & 23.735 \\
  IDR 15T--30T            & 2  &  41.85 & 15  &   9.458 &  1.577 & 0.000 & 11.035 \\
  Above IDR 30T           & 3  & 182.90 & 20  &  83.737 & 10.468 & 0.000 & 94.205 \\
  \midrule
  \rowcolor{lightgray}
  \textbf{Grand Total}    & \textbf{45} & \textbf{373.95} & --- & \textbf{118.234} & \textbf{19.065} & \textbf{0.796} & \textbf{138.095} \\
  \bottomrule
\end{tabular}

\vspace{8pt}
{\footnotesize\color{midgray}
  \textit{Holiday Cost} = investment $\times$ 12\% profit $\times$ CIT forfeited rate $\times$ holiday years. \quad
  \textit{Super Ded. Cost} = investment $\times$ 5\% $\times$ CIT forfeited rate $\times$ 6 years. \quad
  \textit{VAT} = investment $\times$ 35\% $\times$ 11\% (KEK Nongsa facilities only).
}

\vspace{8pt}

% ── Revenue loss: annual + 20-year projection ─────────────────────────────
\section{Government Revenue Foregone: Annual Cost and 20-Year Projection}

\small
\begin{tabular}{lrrrr}
  \toprule
  \textbf{Incentive} &
  \textbf{Annual Cost} &
  \textbf{\% of GDP} &
  \textbf{20-Year Total} &
  \textbf{Notes} \\
  & \textbf{(IDR T/yr)} & \textbf{(2024 GDP)} & \textbf{(IDR T)} & \\
  \midrule
  Tax Holiday (CIT)          & 7.613 & 0.034\% & 118.234 & Avg.\ run-rate; max 20-yr holiday \\
  Super Deduction (PP 78)    & 3.174 & 0.014\% &  19.044 & 6 yrs post-holiday; 5\%/yr deduction \\
  SEZ VAT Exemption (Nongsa) & 0.265 & 0.001\% &   0.796 & One-off; amortised over 3-yr build \\
  \midrule
  \rowcolor{lightgray}
  \textbf{Grand Total}       & \textbf{11.052} & \textbf{0.050\%} & \textbf{138.074} & All incentives combined \\
  \bottomrule
\end{tabular}

\vspace{4pt}
{\footnotesize\color{midgray}
  \textit{Annual cost} for Tax Holiday = sum of (investment $\times$ 12\% $\times$ CIT forfeited rate) across all facilities with an active holiday.
  \textit{Annual cost} for Super Deduction = sum of (investment $\times$ 5\% $\times$ CIT forfeited rate) during the 6-year deduction window.
  \textit{Annual cost} for SEZ VAT = total exemption (IDR 0.796T) $\div$ 3 years assumed construction period.
  Indonesia 2024 nominal GDP = IDR 22,139T (BPS). 20-year totals reflect full contractual periods, not a uniform 20-year stream.
}

\vspace{8pt}

% ── Pillar 2 split ──────────────────────────────────────────────────
\section{Pillar 2 BEPS Impact on Fiscal Cost}

\begin{tabular}{lrrrrrr}
  \toprule
  \textbf{Pillar 2 Status} & \textbf{Facilities} & \textbf{CIT Forfeited} & \textbf{Holiday Cost} & \textbf{Super Ded. Cost} & \textbf{VAT} & \textbf{Total} \\
  \midrule
  In scope (revenue $>$\texteuro{}750M, margin $>$10\%) & 18 & 7\% above 15\% floor  & 14.024 & 2.622 & 0.119 & \textbf{16.765} \\
  Not in scope                                    & 27 & 22\% full statutory    & 104.210 & 16.443 & 0.677 & \textbf{121.330} \\
  \midrule
  \rowcolor{lightgray}
  \textbf{Total}                                  & \textbf{45} & --- & \textbf{118.234} & \textbf{19.065} & \textbf{0.796} & \textbf{138.095} \\
  \bottomrule
\end{tabular}

\vspace{4pt}
{\footnotesize\color{midgray}
  Pillar 2 in-scope entities pay a minimum 15\% effective CIT globally; Indonesia forfeits only 7\% (22\%$-$15\%), as the home jurisdiction collects the top-up tax.
  Non-Pillar 2 entities pay 0\% during the holiday: Indonesia forfeits the full 22\% statutory rate.
}

\vspace{8pt}

% ── Top 5 exposures ─────────────────────────────────────────────────
\section{Top 5 Fiscal Cost Exposures}

\begin{tabular}{llrrrr}
  \toprule
  \textbf{Facility} & \textbf{Operator} & \textbf{P2} & \textbf{Holiday Cost} & \textbf{Super Ded.} & \textbf{Total (IDR T)} \\
  \midrule
  CGK4 Jatiluhur (500 MW)       & BDx Indonesia         & No  & 40.920 & 5.115 & \textbf{46.035} \\
  CGK Campus (500 MW)           & Digital Edge (Indonet)& No  & 36.828 & 4.604 & \textbf{41.432} \\
  JC3 (120 MW)                  & Princeton Digital     & No  &  6.138 & 1.023 & \textbf{7.161}  \\
  Cikarang DC (144 MW)          & DAMAC Digital         & Yes &  5.989 & 0.749 & \textbf{6.738}  \\
  Indonesia Central Cloud Region & Microsoft            & Yes &  3.320 & 0.554 & \textbf{3.874}  \\
  \bottomrule
\end{tabular}

\vspace{4pt}
{\footnotesize\color{midgray}
  BDx and Digital Edge together account for \textbf{IDR 87.5T} (63\%) of the grand total fiscal cost,
  driven by 500 MW capacity commitments on 20-year holidays with no Pillar 2 constraint.
}

\clearpage

% ═══════════════════════════════════════════════════════════════════
%  PAGE 2 — ASSUMPTIONS
% ═══════════════════════════════════════════════════════════════════

{\centering
  {\LARGE\bfseries\color{navy} Methodology \& Key Assumptions}\\[3pt]
  {\small\color{midgray} Indonesia Data Centre Fiscal Cost Analysis --- March 2026}\\[6pt]
}
\noindent\color{navy}\rule{\linewidth}{1.5pt}

\vspace{6pt}

\section{A. Investment Values}
\begin{enumerate}
  \item \textbf{Primary sources:} Confirmed investment figures are drawn from operator press releases, project financing announcements (DBS/UOB, BCA, IIF, green-loan facilities), CEO statements, and due-diligence media (Data Centre Dynamics, Jakarta Globe, Kontan).
  \item \textbf{Imputation for undisclosed values:} Where no public figure exists, investment is estimated using a tiered cost-per-MW model: Tier IV at USD 10M/MW (confirmed from DCI Deputy President benchmark); standard Tier III at USD 6.5M/MW (weighted average across 8 confirmed anchor points); legacy Tier III ($<$2015) at USD 4--5M/MW.
  \item \textbf{Reference investment for tax computations:} For operational facilities, \texttt{investment\_idr\_trillion} is used. For Under Construction and Planned facilities, \texttt{planned\_investment\_idr\_trillion} is used. For K2 JKT1 (both columns populated), the larger full-campus planned value (IDR 5.92T) is applied as it represents the total committed project.
  \item \textbf{Exchange rate:} IDR 15,500 / USD 1 (March 2026 spot mid-rate).
\end{enumerate}

\section{B. Profitability Assumption}
\begin{enumerate}[resume]
  \item \textbf{Annual profit = investment $\times$ 12\%.} This is a sector-level approximation of steady-state EBIT margin on total capital deployed, consistent with published Tier III/IV hyperscale data centre benchmarks (typical range: 10--15\% EBIT on invested capital).
  \item This is a simplification: it ignores ramp-up periods, capital recycling, leverage effects, and asset-light structures (e.g.\ Oracle leasing from DayOne).
\end{enumerate}

\section{C. Tax Holiday (PMK 69/2024)}
\begin{enumerate}[resume]
  \item \textbf{CIT statutory rate:} 22\% (Indonesia corporate income tax, effective from FY2023 onward under UU HPP).
  \item \textbf{Tiered holiday structure} (from PMK 69/2024, amending PMK 130/PMK.010/2020):
  \begin{itemize}
    \item Below IDR 100B: no holiday (below threshold)
    \item IDR 100B--500B: 50\% CIT reduction, 5 years
    \item IDR 500B--1T: 100\% CIT reduction, 5 years
    \item IDR 1T--5T: 100\% reduction, 7 years
    \item IDR 5T--15T: 100\% reduction, 10 years
    \item IDR 15T--30T: 100\% reduction, 15 years
    \item Above IDR 30T: 100\% reduction, 20 years
  \end{itemize}
  \item \textbf{Fiscal cost formula:} Investment $\times$ 12\% (annual profit) $\times$ CIT forfeited rate $\times$ holiday years.
  \item \textbf{Post-holiday transition:} Under PMK 69/2024, a 2-year partial CIT reduction (25--50\%) applies immediately after the holiday ends. This transition period is \textit{not separately costed} in this analysis.
\end{enumerate}

\section{D. BEPS Pillar 2 --- Global Minimum Tax}
\begin{enumerate}[resume]
  \item \textbf{Scope criteria:} An operator is classified as Pillar 2 in-scope if its MNE group meets \textit{both}: (i) consolidated group revenue exceeding \texteuro{}750M, and (ii) a group profitability (operating income / revenue) exceeding 10\%.
  \item \textbf{CIT forfeited (Pillar 2 in scope):} Indonesia forfeits only 7\% (= 22\% statutory $-$ 15\% GloBE minimum floor). The remaining 15\% is collected by the parent jurisdiction as a top-up tax (Qualified Domestic Minimum Top-up Tax or Income Inclusion Rule).
  \item \textbf{CIT forfeited (not in scope):} Indonesia forfeits the full 22\% statutory rate.
  \item \textbf{Result:} 18 of 45 facilities (40\%) are Pillar 2 in-scope, representing IDR 16.8T (12\%) of the total fiscal cost. Pillar 2 meaningfully limits Indonesia's revenue loss only for the in-scope minority.
\end{enumerate}

\section{E. Super Deduction --- Tax Allowance (PP 78/2019)}
\begin{enumerate}[resume]
  \item \textbf{Legal basis:} Government Regulation No.\ 78/2019 (updated by PMK 81/2024) provides a 30\% net income reduction on total realized investment, applied at 5\% per year over 6 years starting after the tax holiday ends.
  \item \textbf{Eligible sectors:} Data centres (KBLI 63111 --- digital hosting and data processing) are explicitly included in the 183 qualifying business fields.
  \item \textbf{Annual deduction:} Investment $\times$ 5\% (reduces taxable income for 6 years post-holiday).
  \item \textbf{Taxable income during super deduction period:} Investment $\times$ (12\% $-$ 5\%) = Investment $\times$ 7\%.
  \item \textbf{Fiscal cost of super deduction:} Annual deduction $\times$ CIT forfeited rate $\times$ 6 years (7\% for Pillar 2 in-scope; 22\% for others).
  \item \textbf{Exception:} Facilities below IDR 100B (no holiday) receive no post-holiday super deduction benefit under this analysis.
\end{enumerate}

\section{F. SEZ VAT Exemption}
\begin{enumerate}[resume]
  \item \textbf{Qualifying SEZ:} Only facilities within KEK Nongsa Digital Park (Batam), designated by Presidential Regulation on 8 June 2021. No other facility in the portfolio has confirmed formal KEK status.
  \item \textbf{VAT rate:} Effective 11\% (12\% statutory $\times$ 11/12 mechanism under PMK-131/2024 for standard capital goods).
  \item \textbf{Exempt investment basis:} 35\% of total investment is assumed to represent VAT-able capital goods (plant and equipment, fit-out). Land, financial instruments, and services subject to reverse-charge are excluded.
  \item \textbf{VAT exemption formula:} Investment $\times$ 35\% $\times$ 11\%.
  \item \textbf{Four SEZ facilities:} PDG Batam Campus (IDR 0.401T), DayOne Nongsa (IDR 0.258T), Oracle Nongsa (IDR 0.119T), Gaw Capital Batam (IDR 0.018T). Total: IDR 0.796T.
\end{enumerate}

\vspace{10pt}
\noindent\footnotesize\color{midgray}
\textit{Disclaimer: This analysis is based on publicly available information, operator disclosures, and sector benchmarks. All investment figures, particularly those marked as imputed (\texttt{investment\_imputed = 1}), carry estimation uncertainty. Fiscal cost figures should be interpreted as indicative order-of-magnitude assessments, not legal or accounting determinations. The interaction between PMK 69/2024 tax holidays, PP 78/2019 tax allowances, and OECD Pillar 2 GloBE rules is subject to Indonesian implementing legislation still under development as of March 2026.}

\clearpage

% ═══════════════════════════════════════════════════════════════════
%  PAGE 3 — FISCAL SCENARIOS
% ═══════════════════════════════════════════════════════════════════

{\centering
  {\LARGE\bfseries\color{navy} Fiscal Scenario Analysis}\\[3pt]
  {\large\color{teal} Scenarios A \& B --- Expanded Incentive Cost Projections}\\[2pt]
  {\small\color{midgray} March 2026 | GDP reference: IDR 22,139T (2024 BPS); 2030 est.\ IDR 34,167T (7.5\% p.a.)}\\[6pt]
}
\noindent\color{navy}\rule{\linewidth}{1.5pt}

\vspace{8pt}

% ── Scenario A ─────────────────────────────────────────────────────
\section{Scenario A --- Expansion to 2,000 MW by 2030}

Scenario A models an expansion from the 526~MW operational baseline to 2,000~MW by 2030, adding 1,474~MW of new capacity at IDR~111B/MW (blended benchmark). New entrants are modelled at 43\% GMT-subject MNEs / 57\% non-GMT, with a 10-year CIT holiday. No new SEZ concessions are assumed.

\small
\begin{tabular}{lrr}
  \toprule
  \textbf{Metric} & \textbf{Baseline (526 MW)} & \textbf{Scenario A --- 2,000 MW (2030)} \\
  \midrule
  Operational MW                      & 526~MW      & 2,000~MW \\
  Incremental MW                      & ---         & +1,474~MW \\
  Total investment (cumulative)       & IDR~76.1T   & IDR~239.7T \\
  Annual CIT cost (peak, 10-yr avg.)  & IDR~1.57T   & IDR~4.62T \\
  Annual cost (\% 2024 GDP)           & 0.007\%     & 0.021\% \\
  Annual cost (\% 2030 GDP est.)      & 0.005\%     & 0.014\% \\
  Cumulative fiscal cost              & IDR~18.0T   & IDR~53.2T \\
  \bottomrule
\end{tabular}

\vspace{4pt}
{\footnotesize\color{midgray}
  Annual CIT cost = cumulative CIT holiday $\div$ 10~years. Scenario~A incremental: CIT~+IDR~30.5T, SD~+IDR~4.7T, total~+IDR~35.2T on top of baseline.
  Future entrants modelled at 43\% GMT (7\% net CIT forfeiture) / 57\% non-GMT (22\% forfeiture); 10-yr holiday assumed for all new builds.
}

\vspace{10pt}

% ── Scenario B ─────────────────────────────────────────────────────
\section{Scenario B --- Expansion to 3,000 MW by 2030}

Scenario B models a full expansion to 3,000~MW by 2030, adding 2,474~MW from the 526~MW baseline at IDR~111B/MW. Entrant mix and holiday structure are identical to Scenario~A.

\small
\begin{tabular}{lrr}
  \toprule
  \textbf{Metric} & \textbf{Baseline (526 MW)} & \textbf{Scenario B --- 3,000 MW (2030)} \\
  \midrule
  Operational MW                      & 526~MW      & 3,000~MW \\
  Incremental MW                      & ---         & +2,474~MW \\
  Total investment (cumulative)       & IDR~76.1T   & IDR~350.7T \\
  Annual CIT cost (peak, 10-yr avg.)  & IDR~1.57T   & IDR~6.69T \\
  Annual cost (\% 2024 GDP)           & 0.007\%     & 0.030\% \\
  Annual cost (\% 2030 GDP est.)      & 0.005\%     & 0.020\% \\
  Cumulative fiscal cost              & IDR~18.0T   & IDR~77.1T \\
  \bottomrule
\end{tabular}

\vspace{4pt}
{\footnotesize\color{midgray}
  Scenario~B incremental (on top of baseline): CIT~+IDR~51.2T, SD~+IDR~7.8T, total~+IDR~59.1T.
  Same 43\%/57\% GMT/non-GMT split and 10-yr holiday as Scenario~A.
}

\vspace{10pt}

% ── VAT/Import Duty ─────────────────────────────────────────────────
\section{VAT / Import Duty --- Universal Extension (Combined Scenarios)}

Extends the existing SEZ VAT exemption (11\%) and import duty waiver (2.5\%) to all operators on the equipment share (35\% of investment). Shown only in combination with Scenario~A or~B; evaluated on the respective investment base.

\small
\begin{tabular}{lrrrr}
  \toprule
  \textbf{Component} & \textbf{Equip.\ Base (IDR T)} & \textbf{VAT (IDR T)} & \textbf{Duty (IDR T)} & \textbf{Total (IDR T)} \\
  \midrule
  Combined A + VAT (2,000 MW) & 83.9 & 9.228 & 2.097 & \textbf{11.326} \\
  Combined B + VAT (3,000 MW) & 122.7 & 13.500 & 3.068 & \textbf{16.569} \\
  \bottomrule
\end{tabular}

\vspace{4pt}
{\footnotesize\color{midgray}
  Equipment base = 35\% of total investment (Scenario~A: IDR~239.7T; Scenario~B: IDR~350.7T). Annualised over 3-year construction phase: Combined~A IDR~3.78T/yr; Combined~B IDR~5.52T/yr.
}

\vspace{10pt}

% ── Combined Summary ───────────────────────────────────────────────
\section{Combined Scenario Summary}

\small
\begin{tabular}{lrrrr}
  \toprule
  \textbf{Scenario} & \textbf{Annual Cost} & \textbf{\% 2024 GDP} & \textbf{\% 2030 GDP} & \textbf{Cumulative (IDR T)} \\
  & \textbf{(IDR T/yr)} & & \textbf{(est.)} & \\
  \midrule
  Baseline (526 MW operational)  & 1.57  & 0.007\% & 0.005\% &  18.0 \\
  \rowcolor{lightgray}
  \textbf{Scenario A} --- 2,000 MW by 2030 & 4.62 & 0.021\% & 0.014\% & 53.2 \\
  \textbf{Scenario B} --- 3,000 MW by 2030 & 6.69 & 0.030\% & 0.020\% & 77.1 \\
  \rowcolor{lightgray}
  \textbf{Combined A + VAT Exemption}      & 8.40 & 0.038\% & 0.025\% & 64.4 \\
  \textbf{Combined B + VAT Exemption}      & 12.22 & 0.055\% & 0.036\% & 93.5 \\
  \bottomrule
\end{tabular}

\vspace{4pt}
{\footnotesize\color{midgray}
  Annual cost = annual peak CIT + annualised VAT+Duty (÷3 yr). Cumulative = total over incentive periods.
  Combined A + VAT: Scenario~A CIT/SD plus IDR~11.3T VAT+Duty on IDR~83.9T equipment base.
  Combined B + VAT: Scenario~B CIT/SD plus IDR~16.6T VAT+Duty on IDR~122.7T equipment base.
  GDP 2024 = IDR 22,139T (BPS); 2030 GDP estimated at IDR 34,167T assuming 7.5\% nominal annual growth.
}

\end{document}
